\section{Original-scale derivative costs of the actual compatibility defects}
\label{mcq:section}

The state and increment in this section are exactly the finite source state
and actual third-growth curl constructed in the preceding reader.
Write \(u_*=\langle u_*\rangle_\theta+w_*\), \(v=w_y\), and
\[
 C=\langle w_*\otimes v+v\otimes w_*+v\otimes v\rangle,
\]
where the outer brackets are the normalized angular and auxiliary
averages before physical evaluation. Every component below is in the
original oriented cylindrical orthonormal basis. The full moment
completion and source pressure calculation prove the formulas used here;
we also derive the scalar compatibility formulas directly to fix all signs.
Throughout this section the unadorned symbols \(P,J_\theta,J_z\)
denote the defect increments \(\Delta P,\Delta J_\theta,\Delta J_z\)
caused by this actual wave addition. They are not the full old-plus-new
source defects. The unchanged old defects must be added when applying
the source correction to the whole updated state.

\subsection{Actual mixed derivatives of the third-growth curl}

Retain the fixed finite set of source labels \(j\), the actual clocks
\(t_{c,j}=t_b+\alpha_jv_j\), and the two full shape fields \(W_{ja}\).
The original data \(t_b,r_a,S_3,u_i\) are independent of these spatial
labels. In particular
\[
 v=\sum_{j,a}Q_j^{-A}y_a(t_{c,j})W_{ja},\quad
 A=\tfrac12+h,\quad D=\tfrac12-h,\quad
 \partial_t t_{c,j}=\alpha_jQ_j^{-1-h},\quad \partial_z t_{c,j}=0 .
\]
Let \(Y_{1,n}=T_n\) and \(Y_{2,n}=O_n\) be the complete finite
derivative arrays already proved for the actual coupled solution.
Thus \(Y_{1,0}=(A_0/\mu)\widehat\lambda_3^{-7/8}\),
\(Y_{2,0}=2A_0\widehat\lambda_3^{1/8}/(\sigma_2\sqrt{n_i})\),
\(Y_{1,1}=2\gamma_iY_{1,0}\), and
\(Y_{2,1}=\gamma_iA_0\widehat\lambda_3^{1/8}/(\sigma_2\sqrt{n_i})\).
Higher entries are those closed recurrences in the retained moving-parent
data; no independent amplitude bound is introduced here.

For a pair \(I=(a,b)\in\mathbb N^2\), put
\(\partial^I=\partial_z^a\partial_t^b\) and
\(I_z=I+(1,0)\). On a fixed chart the full time operator is
\(\mathfrak t_{*,j}=-\varepsilon_j\partial_{T_j}+c_{i,j}N_{i,j}\).
The exact mixed derivative formula is
\begin{equation}\label{mcq:actual_derivative}
\begin{split}
 \partial_z^a\partial_t^b v
 =\sum_{j,c}\sum_{s=0}^b&
 \binom bs Q_j^{-A-aD-b(1+h)}\alpha_j^s
 y_c^{(s)}(t_{c,j})\\
 &\hspace{4mm}\cdot
 \partial_{Z_j}^a\mathfrak t_{*,j}^{\,b-s}W_{jc}.
\end{split}
\end{equation}
All operators on the shapes include their phase, frame and cutoff
derivatives. The shape time operator is the full derivative after
auxiliary evaluation, not a derivative that holds the fast phase fixed.
To prove the identity, differentiate the physical product \(b\) times
in time. If \(s\) of the derivatives act on \(y_c\), its factor is
\((\alpha_jQ_j^{-1-h})^sy_c^{(s)}\); the other \(b-s\) derivatives
give \(Q_j^{-(b-s)(1+h)}\mathfrak t_{*,j}^{b-s}W_{jc}\).
There are \(\binom bs\) such choices. The \(a\) axial derivatives act
only on the shape and give \(Q_j^{-aD}\partial_{Z_j}^a\).
The axial derivative commutes with the stated time operator, since
its coefficients are fixed label constants. Combining the powers proves
\eqref{mcq:actual_derivative} without dropping any derivative.

For each \(e\in\{-1,0,1,2\}\), define the precise weighted norm
\[
 \|f\|_{[e]}^2=\int_0^\infty r^e\langle |f|^2\rangle\,dr,\qquad
 \|F\|_{[e],j}^2=\int_0^\infty R_j^e\langle |F|^2\rangle\,dR_j.
\]
Every field being integrated is supported in a compact annulus away
from zero, so the negative weight introduces no singular integral.
Define the explicit actual-shape budget
\begin{equation}\label{mcq:V}
\begin{split}
 \mathcal V_{e,(a,b)}
 =\sum_{j,c}\sum_{s=0}^b&
 \binom bs \alpha_j^sY_{c,s}
 Q_j^{(e+1)/4-A-aD-b(1+h)}\\
 &\hspace{4mm}\cdot
 \|\partial_{Z_j}^a\mathfrak t_{*,j}^{\,b-s}W_{jc}\|_{[e],j}.
\end{split}
\end{equation}
Then \(\|\partial^I v\|_{[e]}\le\mathcal V_{e,I}\).
Indeed the measure identity
\(r^e\,dr=Q_j^{(e+1)/2}R_j^e\,dR_j\), valid also for \(e=-1\),
and the triangle inequality applied to \eqref{mcq:actual_derivative}
give this bound. There is no inter-label orthogonality assumption.
Put \(\mathcal W_{e,I}=\|\partial^I w_*\|_{[e]}\), restricted to a
fixed radial set enclosing the supports of the actual new fields
and all derivatives at issue. The old wave includes its existing curl
corrections. These are norms of the specified source state, not missing
amplitude variables.

\subsection{A contraction estimate that retains every product derivative}

For \(J=(j_1,j_2)\le I=(a,b)\), set
\(\binom IJ=\binom a{j_1}\binom b{j_2}\), and define
\begin{equation}\label{mcq:B}
\begin{split}
 \mathcal B_{e,I}=\sum_{J\le I}\binom IJ
 \bigl(&\mathcal W_{e,J}\mathcal V_{e,I-J}
       +\mathcal V_{e,J}\mathcal W_{e,I-J}\\
       &+\mathcal V_{e,J}\mathcal V_{e,I-J}\bigr).
\end{split}
\end{equation}
If \(E\) is any fixed real symmetric three-by-three matrix in
cylindrical components, then
\begin{equation}\label{mcq:contraction}
 \left|\partial^I\int r^e(E:C)\,dr\right|
       \le \|E\|_{\mathrm{op}}\mathcal B_{e,I}.
\end{equation}
For a single outer product,
\(E:(f\otimes g)=f^TEg\), whose absolute value is at most
\(\|E\|_{\mathrm{op}}|f||g|\). The repeated product rule gives exactly
the sum over \(J\le I\) in \eqref{mcq:B}. Cauchy--Schwarz on the
original weighted angular-auxiliary-radial measure gives each displayed
product of norms. Compact smooth support justifies differentiation
under the integral. Axial and time derivatives do not differentiate
the cylindrical basis. This proves \eqref{mcq:contraction} for every
specified order.
For the time derivative, the full lifted physical operator is
\(\mathcal D_t=\partial_t+N_{\rm abs}\), with \(N_{\rm abs}\) a
constant torus direction. The Haar integral of \(N_{\rm abs}f\) is
zero by periodicity. Therefore
\(\partial_t\langle f\rangle=\langle\mathcal D_tf\rangle\).
The same identity applies repeatedly, and axial differentiation commutes
with this average and this constant torus direction. Consequently the
ordinary mixed derivatives of the averaged tensor in
\eqref{mcq:contraction} are exactly the averages of the full physical
derivatives estimated in \eqref{mcq:actual_derivative}.

\subsection{The actual three defects and the pressure-corrected axial flux}

Since the wave increment has zero angular mean and is divergence-free,
the change of the radial mean pressure source is
\[
 g=-(\partial_r+r^{-1})C_{rr}-\partial_zC_{zr}+r^{-1}C_{\theta\theta}.
\]
The three actual source compatibility-defect increments are therefore
\begin{equation}\label{mcq:defects}
\begin{split}
 P&=\int\frac{C_{\theta\theta}-C_{rr}}r\,dr
                    -\partial_z\int C_{zr}\,dr,\\
 J_\theta&=\int r^2C_{z\theta}\,dr,\\
 J_z&=\int r\left(C_{zz}-\tfrac12C_{rr}
                              -\tfrac12C_{\theta\theta}\right)\,dr
                  +\tfrac12\partial_z\int r^2C_{zr}\,dr .
\end{split}
\end{equation}
The first formula is \(\int g\,dr\), using the zero boundary integral
of \(\partial_r C_{rr}\). For the last, integration by parts gives
\[
 \int r^2 g\,dr
  =\int r(C_{rr}+C_{\theta\theta})\,dr
                      -\partial_z\int r^2C_{zr}\,dr .
\]
Substitute it into the source definition
\(J_z=\int rC_{zz}\,dr-\frac12\int r^2g\,dr\).
This proves all three identities in \eqref{mcq:defects}.
In particular the axial source defect is not replaced by the raw flux
\(\int rC_{zz}\,dr\).

Let the source pressure bump be the original
\[
 \rho(r,z,t)=q^{-1/2}\widehat\rho(r/\sqrt q),\qquad
 \int\widehat\rho(x)\,dx=1,\qquad
 c_0=\tfrac12\int x^2\widehat\rho(x)\,dx .
\]
It has \(\int\rho\,dr=1\) and
\(c_\rho=\frac12\int r^2\rho\,dr=c_0q\).
The auxiliary average of the compact source pressure increment is
\[
 \overline{\Delta p_m}(r)=\int_0^r(g(s)-\rho(s)P)\,ds .
\]
Its derivative is \(g-\rho P\), and its integrand has zero total
radial integral, giving compact support. Its weighted pressure moment is
\(\int r\overline{\Delta p_m}\,dr=-\frac12\int r^2g\,dr+c_\rho P\).
Here is the exact map to the full auxiliary-dependent source operation.
Let \(\Delta g_r\) be the angularly averaged source increment before
auxiliary averaging, so that \(\langle\Delta g_r\rangle_Y=g\), and
put \(f=\Delta g_r-\rho P\). The source's shifted primitive uses
\[
 \begin{split}
 \mathcal I f(r,Y)&=\int_0^r
       f(s,Y+(s^{d_r}-r^{d_r})v_r)\,ds,\\
 \mathcal J f(r,Y)&=\int_0^\infty
       f(s,Y+(s^{d_r}-r^{d_r})v_r)\,ds,\qquad
 \mathcal I_cf=\mathcal If-\chi_m\mathcal Jf .
 \end{split}
\]
The fixed source exponent and direction are the original \(d_r,v_r\),
and \(\chi_m\) is its inner-zero, outer-one radial cutoff. The full
pressure increment is \(\Delta p_m=\mathcal I_cf\).
Haar translation invariance gives
\(\langle\mathcal Jf\rangle_Y=\int(g-\rho P)\,dr=0\) and
\(\langle\mathcal If\rangle_Y=\int_0^r(g-\rho P)\,ds\).
This proves the asserted averaged pressure formula from the full map.
The lifted radial derivative cancels the shift derivative and yields
\[
 \mathcal D_r\Delta p_m-\Delta g_r
       =-\rho P-(\partial_r\chi_m)\mathcal Jf .
\]
The second term has zero auxiliary average but need not vanish before
physical evaluation. No estimate for its singular-time flatness is
inferred from its zero average.
Consequently the double-averaged axial residual-moment increment after
this reconstruction is
\[
 \partial_zK_z,\qquad K_z=J_z+c_0qP .
\]
The full product, including the derivatives of \(q\), is retained.

The following simultaneous quantitative estimates hold at every finite
mixed order:
\begin{equation}\label{mcq:defect_bounds}
\begin{split}
 |\partial^IP|&\le
       \mathcal B_{-1,I}+\tfrac12\mathcal B_{0,I_z},\\
 |\partial^IJ_\theta|&\le\tfrac12\mathcal B_{2,I},\\
 |\partial^IJ_z|&\le
       \mathcal B_{1,I}+\tfrac14\mathcal B_{2,I_z},\\
 |\partial^IK_z|&\le
       \mathcal B_{1,I}+\tfrac14\mathcal B_{2,I_z}\\
 &\quad+|c_0|\sum_{J\le I}\binom IJ
       |\partial^Jq|\left(
       \mathcal B_{-1,I-J}+\tfrac12\mathcal B_{0,(I-J)_z}\right).
\end{split}
\end{equation}
For the first line the diagonal contraction matrix is
\(\operatorname{diag}(-1,1,0)\), of norm one, whereas the \(rz\)
entry is the contraction against the symmetric matrix with \(rz,zr\)
entries \(1/2\), of norm \(1/2\). For the second line use the
analogous \(z\theta\) matrix. For the third the diagonal matrix is
\(\operatorname{diag}(-1/2,-1/2,1)\), of norm one, and the
coefficient \(1/2\) multiplying \(C_{zr}\) contributes an additional
factor \(1/2\). These facts and \eqref{mcq:contraction} prove the first
three bounds. The exact product rule for \(qP\) proves the fourth.
Replacing \(I\) by \(I_z\) gives bounds for the two residual moments
that actually enter the source radial inverse. Thus every pressure,
radial and axial coupling has a bound in the original shape fields
and third-growth constants.

\subsection{Exact derivatives of the original moving scale and bumps}

The source scale is determined by
\[
 q-z^2q^{2h}=1-t,\qquad z=q^D\eta,\qquad
 D=\tfrac12-h,\qquad 0<h<1/100,\quad -1<\eta<1 .
\]
Set \(B(\eta)=1-2h\eta^2>1-2h>0\). Implicit differentiation of the
first equation, followed by the displayed formula for \(z\), gives
\begin{equation}\label{mcq:q_derivatives}
\begin{gathered}
 q_z=\frac{2\eta q^A}{B},\qquad q_t=-\frac1B,\\
 \eta_z=\frac{q^{-D}(1-\eta^2)}B,\qquad
 \eta_t=\frac{D\eta q^{-1}}B .
\end{gathered}
\end{equation}
For example \(q_z=2zq^{2h}/B\), and \(D+2h=A\).
Differentiating \(\eta=zq^{-D}\) gives the last two formulas;
the identity \(2D+2h=1\) gives the factor \(1-\eta^2\).
This proves the derivative dictionary without assigning a sign to \(q_z\).

Put \(x=r/\sqrt q\). For a smooth function \(f(x,\eta)\) define the
explicit differential expressions
\[
 \mathcal Z_\sigma f=
 \frac{2\sigma\eta f-\eta x f_x+(1-\eta^2)f_\eta}{B},\qquad
 \mathcal T_\sigma f=
 \frac{-\sigma f+\frac12x f_x+D\eta f_\eta}{B}.
\]
Equations \eqref{mcq:q_derivatives} and the ordinary chain rule prove
\begin{equation}\label{mcq:dilation_dictionary}
 \partial_z(q^\sigma f)=q^{\sigma-D}\mathcal Z_\sigma f,\quad
 \partial_t(q^\sigma f)=q^{\sigma-1}\mathcal T_\sigma f,\quad
 \partial_r(q^\sigma f)=q^{\sigma-1/2}f_x .
\end{equation}
For the last formula \(q,\eta\) are independent of \(r\).
For the first two, differentiation of \(x\) gives respectively
\(-\eta xq^{-D}/B\) and \(xq^{-1}/(2B)\), proving every term.

In particular let the original radial moment bump be
\[
 b_e=q^{-(e+1)/2}\widehat b_e(x),\qquad
 \int x^e\widehat b_e(x)\,dx=1 .
\]
For any \(k,a,b\ge0\), its derivative is exactly
\begin{equation}\label{mcq:bump_derivatives}
 \partial_r^k\partial_z^a\partial_t^b b_e
 =q^{-(e+1+k)/2-aD-b} f_{e;k,a,b}(x,\eta).
\end{equation}
Here the coefficient is given by a finite recursion: start at
\(\sigma=-(e+1+k)/2\) and \(f=\widehat b_e^{(k)}\);
repeat \(f\leftarrow\mathcal Z_\sigma f\),
\(\sigma\leftarrow\sigma-D\), \(a\) times, and then
\(f\leftarrow\mathcal T_\sigma f\),
\(\sigma\leftarrow\sigma-1\), \(b\) times.
The result is \(f_{e;k,a,b}\). This proves the identity by induction
using \eqref{mcq:dilation_dictionary}. Smooth mixed derivatives commute,
so this fixed order computes the stated derivative.
Every coefficient is explicitly a finite rational expression in \(x,\eta\),
the retained constants, and derivatives of the specified bump, with
nonzero denominator a power of \(B\). On its compact \(x\) support
and \(|\eta|\le1\) its maximum is finite and computable directly from
that expression. Thus \eqref{mcq:bump_derivatives} supplies the original
power of \(q\) at every order.
The pressure bump \(\rho\) uses exactly this recursion with \(e=0\).
Likewise every derivative in \eqref{mcq:defect_bounds} is explicit:
\(\partial_z^a\partial_t^bq=q^{1-aD-b}Q_{a,b}(\eta)\), where the same
recursion starts with \(\sigma=1,f=1\).

For completeness, the weighted cumulative bump has the exact form
\[
 \beta_e(r,z,t)=\int_0^rs^eb_e(s,z,t)\,ds
              =\int_0^{r/\sqrt q}x^e\widehat b_e(x)\,dx .
\]
Its axial derivative is
\begin{equation}\label{mcq:bump_movement}
 \partial_z\beta_e
 =-\frac{q_z}{2q}\,r^{e+1}b_e .
\end{equation}
This follows by differentiating its upper endpoint, which is
\(-rq_z/(2q^{3/2})\). Inserting the definition of \(b_e\) gives
the displayed original physical factors. Therefore replacing the
primitive of \(F_e-b_e\partial_zJ_e\) by the exact primitive of
\(F_e-\partial_z(b_eJ_e)\) changes the stress coefficient by
\[
 -\frac{rq_z}{2q}\,b_eJ_e .
\]
This is the moving-bump term required in the complete stress construction.
Its sign and all further derivatives are fixed by
\eqref{mcq:q_derivatives}--\eqref{mcq:bump_movement}; it is not set to zero
by holding the physical source scale fixed during axial differentiation.

These calculations supply finite quantitative inputs for the actual
source pressure and moment operations. They do not identify the
distinct direct candidate with a completed infinite source correction
sequence: the complete source-operation formulas retain the new
velocity, covariance, pressure and nonlinear defects that such a step
actually creates.
